\documentclass{article}
\usepackage{amsmath, amssymb, amsthm}
\usepackage{hyperref}
\usepackage{geometry}
\geometry{margin=1in}

\title{Erd\H{o}s Problem \#712}
\date{Last updated: 2025-08-31}
\author{Source: erdosproblems.com}

\begin{document}
\maketitle

\noindent\textbf{Status:} OPEN \quad \textbf{Prize: \$500}

\noindent\textbf{Tags:} graph theory, turan number, hypergraphs

\medskip

\noindent\textbf{Problem Statement:}

Determine, for any $k&#62;r&#62;2$, the value of\[\frac{\mathrm{ex}_r(n,K_k^r)}{\binom{n}{r}},\]where $\mathrm{ex}_r(n,K_k^r)$ is the largest number of $r$-edges which can placed on $n$ vertices so that there exists no set of $k$ vertices which is covered by all $\binom{k}{r}$ possible $r$-edges.




Tur\'{a n proved} that, when $r=2$, this limit is\[\frac{1}{2}\left(1-\frac{1}{k-1}\right).\]Erd\H{o}s \cite{Er81} offered \$500 for the determination of this value for any fixed $k&gt;r&gt;2$, and \$1000 for 'clearing up the whole set of problems'. See also [500] for the case $r=3$ and $k=4$.




References

[Er81] Erd\H{o}s, P., On the combinatorial problems which I would most like to see solved . Combinatorica (1981), 25-42.

\noindent\textbf{Related OEIS sequences:} possible


\bigskip
\noindent\small{Source: \url{https://www.erdosproblems.com/712}}
\end{document}
